\chapter{Conclusion and Future Work}
\label{ch:conclusion}

\section{Conclusions}

This thesis presented the design, implementation, and evaluation of an automated label
validation pipeline for IKEA of Sweden's communication label catalogue. The system addresses
three research questions posed in the introduction.

\textbf{RQ1 — Classification accuracy}: The system achieves a validation F1 of~0.865
(precision~88.3\%, recall~84.7\%) in distinguishing compliant from non-compliant labels,
using a combination of multi-engine \ac{OCR}, GS1 barcode decoding, and rule-based
field checking. The accuracy is bounded primarily by the completeness of the rule corpus
rather than by the quality of the extraction—two of the four total errors in the test set
involved rules that are not yet encoded.

\textbf{RQ2 — Explainability}: Every validation decision is fully traceable to a specific
rule, with extracted vs.\ expected values presented for each failing check. Domain expert
evaluation confirmed that the output is clear, actionable, and sufficient for a QA operator
to locate the violation in the specification document within 30 seconds. The rule-based
architecture provides this interpretability by construction, without requiring post-hoc
explanation methods.

\textbf{RQ3 — Contributing features}: Label-type identification benefits most from
barcode type presence (DataMatrix vs.\ ITF-only schemas are highly distinguishable). Field
extraction accuracy is highest for structured identifiers with fixed patterns (supplier name,
article number) and lowest for free-format fields (dimensions, weights). Validation accuracy
is most sensitive to the completeness and correctness of the schema rules rather than to
OCR accuracy—implying that schema curation is the highest-leverage investment for
improving the system.

Beyond the research questions, the project delivers a concrete industrial artefact: a
fully functional validation pipeline accessible via a \ac{REST} \ac{API} and an interactive
Streamlit \ac{UI}, covering all 34 label types in IKEA of Sweden's current catalogue,
with a ground-truth rule corpus of~342 field declarations and~107 validation rules.

\section{Future Work}

Several directions are identified for extending this work.

\subsection{Layout and Spatial Validation}

The \texttt{layout\_checker.py} module is currently a stub. Implementing zone-order
checking (verifying that zones appear in the canonical left-to-right sequence) and
element-in-zone checking (verifying that each field is positioned in its expected zone)
would address the two false-negative cases in the current evaluation and likely improve
recall substantially.

\subsection{Conditional Rules}

The schema format should be extended to support conditional field presence rules of the
form: ``field $X$ is required if and only if field $Y$ is also present.'' This would allow
encoding the date-zone activation rule (YYWW stamp required only when date section is
activated) and several other conditional constraints that are currently approximated as
optional fields.

\subsection{Neural Field Extraction}

For low-quality raster inputs (scans, photographs), the regex-based field mapper performs
poorly when OCR output is noisy. A LayoutLMv3 model fine-tuned on the IKEA label
dataset—using the~238 corrected annotations as training data—could provide substantially
better field extraction on difficult inputs. The \texttt{/api/export-jsonl} endpoint was
designed specifically to facilitate this future fine-tuning step.

\subsection{PaddleOCR Integration}

PaddleOCR (PP-OCRv5) was the highest-accuracy engine in preliminary benchmarking
(92.3\% character accuracy) but could not be included in the primary evaluation due to
platform incompatibilities. A Linux deployment should replace EasyOCR with PaddleOCR
for a further improvement in ensemble accuracy.

\subsection{Physical Dimension Inference}

If a known physical reference is available (e.g.\ a calibration pattern printed alongside
the label, or the known x-dimension of a decoded barcode), pixel-to-mm conversion becomes
possible. This would enable:
\begin{itemize}
  \item Intra-family label-type disambiguation based on physical label size.
  \item Barcode module-size validation against the ISO/IEC barcode specification.
  \item Font size validation for the visual hierarchy rules.
\end{itemize}

\subsection{Dataset Expansion}

Expanding the dataset to cover all 34 label types and increasing the per-type sample
count to at least~20 would enable statistically meaningful per-type performance analysis
and support the training of neural extraction models. IKEA's \ac{QA} team collects
approximately~50 failed labels per month; integrating this production data stream would
grow the dataset organically over time.

\subsection{API Dashboard and Integration}

A production-grade deployment of the system would require authentication, audit logging,
rate limiting, and integration with IKEA's product information management (\ac{PIM})
system so that the expected article number, origin, and DataMatrix content can be retrieved
automatically from the product database rather than inferred from the label itself. This
would allow the system to catch cases where the label \emph{internally consistent} but
disagrees with the \ac{PIM} record—the most dangerous class of label error.

\section{Final Remarks}

This work demonstrates that a carefully designed rule-based pipeline—augmented with
multi-engine \ac{OCR} and GS1 barcode parsing—can automate the majority of IKEA's
label \ac{QA} workload with high accuracy and full interpretability, using no neural
network training and no GPU hardware. The system is ready for a pilot deployment, and the
architecture is extensible to accommodate the improvements outlined above as the project
matures from a bachelor's thesis prototype to a production service.
